%% OptoNet --- SSRN-ready preprint source.
%% Compile with:  tectonic main.tex   (or pdflatex main.tex)
\documentclass[11pt,a4paper]{article}

\usepackage[T1]{fontenc}
\usepackage[utf8]{inputenc}
\usepackage{geometry}
\geometry{margin=1in}
\usepackage{amsmath,amssymb,bm}
\usepackage{graphicx}
\usepackage{booktabs}
\usepackage{multirow}
\usepackage{array}
\usepackage{xcolor}
\usepackage[hidelinks]{hyperref}
\usepackage{microtype}
\usepackage{enumitem}
\usepackage{caption}
\usepackage{natbib}

% BEGIN AUTO NUMBERS
% Auto-generated by scripts/11_paper_numbers.py -- do not edit
\newcommand{\nsource}{OptoNet simulator benchmark}
\newcommand{\nTrials}{61000}
\newcommand{\nCells}{900}
\newcommand{\nSpikes}{736297}
\newcommand{\nFamilies}{9}
\newcommand{\nTrainFamilies}{5}
\newcommand{\nOodFamilies}{4}
\newcommand{\nSeeds}{3}
\newcommand{\conformalAlpha}{0.100}
\newcommand{\RtwoGlmId}{0.688}
\newcommand{\auPrcGlmId}{0.037}
\newcommand{\auRocGlmId}{0.849}
\newcommand{\maeGlmId}{4.074}
\newcommand{\psthGlmId}{0.277}
\newcommand{\RtwoGlmProto}{-4.155}
\newcommand{\auPrcGlmProto}{0.063}
\newcommand{\RtwoGlmFreq}{0.563}
\newcommand{\RtwoGlmInt}{-0.263}
\newcommand{\psthGlmProto}{0.703}
\newcommand{\nParamsGlm}{1552}
\newcommand{\RtwoCnnId}{0.982}
\newcommand{\auPrcCnnId}{0.045}
\newcommand{\auRocCnnId}{0.870}
\newcommand{\maeCnnId}{1.074}
\newcommand{\psthCnnId}{0.318}
\newcommand{\RtwoCnnProto}{0.963}
\newcommand{\auPrcCnnProto}{0.074}
\newcommand{\RtwoCnnFreq}{0.972}
\newcommand{\RtwoCnnInt}{0.974}
\newcommand{\psthCnnProto}{0.723}
\newcommand{\nParamsCnn}{43921}
\newcommand{\RtwoLstmId}{0.936}
\newcommand{\auPrcLstmId}{0.034}
\newcommand{\auRocLstmId}{0.856}
\newcommand{\maeLstmId}{1.814}
\newcommand{\psthLstmId}{0.197}
\newcommand{\RtwoLstmProto}{0.915}
\newcommand{\auPrcLstmProto}{0.062}
\newcommand{\RtwoLstmFreq}{0.851}
\newcommand{\RtwoLstmInt}{0.837}
\newcommand{\psthLstmProto}{0.684}
\newcommand{\nParamsLstm}{57313}
\newcommand{\RtwoTransformerId}{0.961}
\newcommand{\auPrcTransformerId}{0.041}
\newcommand{\auRocTransformerId}{0.860}
\newcommand{\maeTransformerId}{1.503}
\newcommand{\psthTransformerId}{0.268}
\newcommand{\RtwoTransformerProto}{0.898}
\newcommand{\auPrcTransformerProto}{0.067}
\newcommand{\RtwoTransformerFreq}{0.954}
\newcommand{\RtwoTransformerInt}{0.958}
\newcommand{\psthTransformerProto}{0.694}
\newcommand{\nParamsTransformer}{148672}
\newcommand{\ciGlmLo}{0.640}
\newcommand{\ciGlmHi}{0.734}
\newcommand{\seedSdGlm}{0.011}
\newcommand{\ciCnnLo}{0.980}
\newcommand{\ciCnnHi}{0.983}
\newcommand{\seedSdCnn}{0.000}
\newcommand{\ciLstmLo}{0.928}
\newcommand{\ciLstmHi}{0.943}
\newcommand{\seedSdLstm}{0.006}
\newcommand{\ciTransformerLo}{0.959}
\newcommand{\ciTransformerHi}{0.963}
\newcommand{\seedSdTransformer}{0.005}
\newcommand{\RtwoHgbId}{0.788}
\newcommand{\maeHgbId}{3.853}
\newcommand{\RtwoHgbProto}{-0.462}
\newcommand{\maeHgbProto}{11.008}
\newcommand{\RtwoHgbFreq}{0.792}
\newcommand{\maeHgbFreq}{5.767}
\newcommand{\RtwoHgbInt}{0.780}
\newcommand{\maeHgbInt}{5.337}
\newcommand{\RtwoHgbVal}{0.744}
\newcommand{\maeHgbVal}{3.941}
\newcommand{\auRocHgbDetId}{0.958}
\newcommand{\auPrcHgbDetId}{0.978}
\newcommand{\auRocHgbDetProto}{0.906}
\newcommand{\auPrcHgbDetProto}{0.991}
\newcommand{\covGlmId}{0.898}
\newcommand{\sdGlmId}{0.611}
\newcommand{\covGlmProto}{0.280}
\newcommand{\sdGlmProto}{1.579}
\newcommand{\covGlmFreq}{0.784}
\newcommand{\sdGlmFreq}{1.487}
\newcommand{\covGlmInt}{0.770}
\newcommand{\sdGlmInt}{2.358}
\newcommand{\covCnnId}{0.903}
\newcommand{\sdCnnId}{0.236}
\newcommand{\covCnnProto}{0.780}
\newcommand{\sdCnnProto}{0.396}
\newcommand{\covCnnFreq}{0.802}
\newcommand{\sdCnnFreq}{0.462}
\newcommand{\covCnnInt}{0.832}
\newcommand{\sdCnnInt}{0.384}
\newcommand{\covLstmId}{0.898}
\newcommand{\sdLstmId}{0.388}
\newcommand{\covLstmProto}{0.819}
\newcommand{\sdLstmProto}{0.821}
\newcommand{\covLstmFreq}{0.730}
\newcommand{\sdLstmFreq}{0.727}
\newcommand{\covLstmInt}{0.755}
\newcommand{\sdLstmInt}{0.785}
\newcommand{\covTransformerId}{0.871}
\newcommand{\sdTransformerId}{0.436}
\newcommand{\covTransformerProto}{0.693}
\newcommand{\sdTransformerProto}{0.831}
\newcommand{\covTransformerFreq}{0.783}
\newcommand{\sdTransformerFreq}{1.064}
\newcommand{\covTransformerInt}{0.798}
\newcommand{\sdTransformerInt}{0.823}
\newcommand{\trOracleRtwo}{0.974}
\newcommand{\trOracleAuprc}{0.030}
\newcommand{\trZeroShotRtwo}{0.509}
\newcommand{\trZeroShotAuprc}{0.027}
\newcommand{\trScratchNRtwo}{0.930}
\newcommand{\trScratchNAuprc}{0.028}
\newcommand{\trFewShotNRtwo}{0.966}
\newcommand{\trFewShotNAuprc}{0.030}
\newcommand{\trFinetuneOracleNRtwo}{0.970}
\newcommand{\trFinetuneOracleNAuprc}{0.030}
\newcommand{\holdoutType}{SOM}
\newcommand{\salEnrichGlm}{1.141}
\newcommand{\shuffleAuPrcGlm}{0.0074}
\newcommand{\salEarlyGlm}{1.01}
\newcommand{\salEnrichCnn}{1.506}
\newcommand{\shuffleAuPrcCnn}{0.0076}
\newcommand{\salEarlyCnn}{1.13}
\newcommand{\salEnrichLstm}{0.881}
\newcommand{\shuffleAuPrcLstm}{0.0073}
\newcommand{\salEarlyLstm}{0.99}
\newcommand{\salEnrichTransformer}{0.924}
\newcommand{\shuffleAuPrcTransformer}{0.0076}
\newcommand{\salEarlyTransformer}{2.57}
\newcommand{\attnLocal}{0.287}
\newcommand{\attnLocalUniform}{0.424}
\newcommand{\glmPeakPc}{5}
\newcommand{\glmPeakPv}{4}
\newcommand{\glmPeakSom}{7}
\newcommand{\optErrOurs}{9.042}
\newcommand{\optDoseOurs}{146.436}
\newcommand{\optWithinOurs}{0.146}
\newcommand{\optErrConst}{10.156}
\newcommand{\optDoseConst}{170.588}
\newcommand{\optWithinConst}{0.167}
\newcommand{\optErrRand}{8.323}
\newcommand{\optDoseRand}{148.002}
\newcommand{\optWithinRand}{0.167}
\newcommand{\optErrOracle}{0.646}
\newcommand{\optDoseOracle}{226.667}
\newcommand{\optWithinOracle}{0.156}
\newcommand{\optTarget}{20}
% END AUTO NUMBERS

\definecolor{optoblue}{RGB}{43,108,176}
\definecolor{optoorange}{RGB}{192,86,33}
\newcommand{\opto}{\textsc{OptoNet}}
\newcommand{\former}{\textsc{OptoFormer}}
\newcommand{\Rtwo}{\ensuremath{R^{2}}}
\newcommand{\paperfig}[3]{%
  \begin{figure}[t]\centering
  \IfFileExists{#1}{\includegraphics[width=\linewidth]{#1}}{}
  \caption{#2}\label{#3}\end{figure}}

\title{\bfseries \opto{}: Deep Learning Predicts and Designs\\
Neuron-Type-Specific Responses to Optogenetic Stimulation}
\author{Basit Ali\thanks{Abdul Wali Khan University, Mardan, Pakistan.
Corresponding author: \texttt{basitali.262713@gmail.com}}}
\date{}

\begin{document}
\maketitle

\begin{abstract}
\noindent
Optogenetic experiments are constrained by a hard bottleneck: every candidate
stimulation protocol must be prepared, delivered and recorded before one knows
what a neuron will do. We ask whether this loop can be shortened. We introduce
\opto{}, a reproducible benchmark and model suite for \emph{predicting} and
\emph{designing} the spike train that a light protocol elicits. The benchmark
consists of \nTrials{} trials generated by a biophysically grounded forward
model: adaptive exponential integrate-and-fire neurons of three cell types
(pyramidal, PV-like fast-spiking, SOM-like adapting) driven by a two-state
ChR2 photocurrent with light-dependent activation, slow desensitisation and
voltage-dependent deactivation, plus Ornstein--Uhlenbeck background
fluctuations. Each neuron is drawn from a heterogeneous parameter population,
so no two cells are identical. Nine stimulation families (regular trains, step
pulses, gamma-distributed random trains, doublets, bursts, chirps, sweeps,
sparse stimuli and high-frequency sweeps) are split into a training
distribution and three
out-of-distribution (OOD) test regimes: entirely unseen protocol families,
unseen frequencies ($50$--$120$~Hz) and unseen intensities.

We train and compare four sequence models at 1~kHz---the classical
cell-type-specific linear--nonlinear encoding model (GLM), a causal dilated CNN,
a convolutional LSTM, and \former{}, a causal transformer with FiLM cell
conditioning---against gradient-boosted trees on handcrafted features. In
distribution, deep models explain $\Rtwo \approx \RtwoTransformerId$ of the
variance in evoked spike counts and reach AUPRC $=\auPrcTransformerId$ for
millisecond-level spike detection (base rate $0.008$). Out of distribution the
models separate sharply:
gradient-boosted trees extrapolate reasonably to unseen frequencies and
intensities ($\RtwoHgbFreq$, $\RtwoHgbInt$) but \emph{collapse} on unseen
protocol families ($\RtwoHgbProto$), whereas the sequence models degrade
gracefully ($\RtwoTransformerProto$). We further show that (i) deep ensembles
plus split-conformal prediction flag distribution shift through inflated
epistemic uncertainty and degraded coverage, (ii) a model pretrained on two cell
types can be adapted to a held-out third type with a few dozen labelled cells,
and (iii) most importantly, \emph{the learned model can be run in the loop}:
a surrogate model is used to search for pulse protocols that hit a target firing
output, and the designed protocols are then validated against the
ground-truth simulator, converting a slow empirical design loop into an
in-silico one while measuring how much surrogate bias remains.
Interpretability analyses confirm that the models read the stimulus: gradient
saliency concentrates on the light trace with a strong onset transient, the
learned GLM kernel recovers the fast photocurrent onset used by the forward
model, and shuffling the stimulus across trials destroys performance.
All code, data generators, and results
tables are released.
\end{abstract}

\noindent\textbf{Keywords:} optogenetics, channelrhodopsin, computational
neuroscience, neural encoding models, deep learning, spiking neurons,
distribution shift, predictive uncertainty, closed-loop stimulation design

\noindent\textbf{JEL classification:} C45, C63, I23, Q05

\bigskip

\section{Introduction}

Optogenetics made causal manipulation of neuronal activity routine. A single
transgene, a patterned light stimulus, and an electrode: the causal arrow runs
from light to spikes. That power, however, comes with a severe experimental
bottleneck. The space of possible stimulation protocols is enormous, and the
cost of evaluating any single protocol---surgery, recovery, hours of
recording---is enormous as well. Experiments are therefore designed around
protocols whose effects are qualitatively anticipated from prior art: single
pulses to check expression, trains at a handful of frequencies, and chirps to
inspect resonance. Nothing in the standard workflow tells the experimenter, for
a given cell, which protocol will produce a desired firing pattern.

This is a prediction problem, and machine learning is well suited to it. Given
the time course of a light stimulus and a description of the target neuron,
predict the spike train. If that prediction is accurate, three capabilities
follow immediately. (i)~Protocol \emph{screening}: thousands of candidate
protocols can be evaluated in seconds before touching the rig. (ii)~Protocol
\emph{design}: the model can be inverted---optimised against---to find protocols
that produce a desired output. (iii)~Rapid cell-type \emph{adaptation}: a model
pretrained on one population can be fine-tuned on a handful of cells of a new
type instead of retrained from scratch.

The neural coding literature already contains a well-established answer to the
prediction half of this problem: the generalised linear model (GLM) of
\citep{pillow2008spatiotemporal}, which explains spike trains as a history of
weighted stimulus filters passed through a static non-linearity. GLMs are
interpretable, have strong theoretical grounding, and are the tool of choice for
first-pass analysis. They are, however, shallow by construction: a single
linear-temporal filter followed by a pointwise non-linearity cannot represent
state-dependent phenomena such as spike-frequency adaptation, opsin
desensitisation across a train, or cell-type-specific thresholds. In
optogenetics these effects are not nuisances; they are the reason that
stimulation at 20~Hz and 80~Hz produces qualitatively different behaviour.

Recent deep architectures---dilated convolutional networks, recurrent networks
and transformers---remove the architectural bottleneck. Whether they actually
help for this problem, however, is rarely established on a benchmark designed
to make the answer unambiguous, because published datasets are small, biased
towards a handful of protocols, and lack held-out stimulation regimes.

This paper contributes three things.

\begin{enumerate}[leftmargin=*,itemsep=2pt]
\item \textbf{A benchmark that can give a wrong answer.} \opto{} generates
      \nTrials{} trials from \nCells{} heterogeneous virtual neurons across
      \nFamilies{} stimulation families, with a training distribution and three
      deliberately adversarial test regimes: protocol families never seen in
      training, frequencies never seen in training, and intensities beyond the
      training range. A model that has merely memorised stimulus families
      cannot score well, and the tree baseline demonstrates exactly that failure
      mode.
\item \textbf{A controlled model comparison at spike-time resolution.} Four
      causal sequence models spanning ~1.5k to ~150k parameters, all trained
      with an identical loss, data split and optimiser, and evaluated with
      identical metrics at 1~kHz: count accuracy, spike-detection AUPRC, and
      PSTH correlation.
\item \textbf{Using the model in the loop.} We close the loop with a
      surrogate-driven search over pulse protocols, validate the designed
      protocols against the ground-truth simulator, and quantify both what a
      flexible protocol space buys over a constant-frequency search and how
      much design bias the surrogate still carries.
\end{enumerate}

\noindent
A deliberate scope statement. This is a \emph{simulation} study. The forward
model is not a substitute for real recordings; it is an instrument. Its
advantage is that the ground truth is exactly known, so every number reported
here is a statement about the learning problem rather than about experimental
drift, and every evaluation---including the protocol-design experiment---can be
scored against an oracle. The forward model is built from published mechanisms
\cite{brette2005adex,nagel2003light,foutz2012basis,herman2019minimal} and is
calibrated to reproduce the qualitative phenomenology reported for
channelrhodopsin-driven cortical neurons (Section~\ref{sec:sim}). We discuss
what would and would not transfer to real data in
Section~\ref{sec:limits}.

\section{Related work}

\textbf{Neural encoding models.} The GLM framework
\citep{pillow2008spatiotemporal} and its convolutional and adaptive-filter
variants remain the dominant tools for describing stimulus--spike
relationships, because a learned filter can be read as a receptive field and a
learned coupling can be read as interactions \cite{pillow2008spatiotemporal}.
Adaptation and refractory effects enter through the history term. What GLMs do
not capture is non-stationarity in the \emph{stimulus} channel: the mapping from
light to spikes changes within a trial as opsins desensitise and as the cell's
adaptation variable winds up.

\textbf{Biophysical forward models.} AdEx neurons
\citep{brette2005adex} are a standard compromise between biological richness and
simulability and are used here as the substrate. Detailed photocurrent models
for channelrhodopsin \cite{nagel2003light,foutz2012basis,herman2019minimal}
supply light-dependent activation, desensitisation and voltage-dependent
deactivation.

\textbf{Deep sequence models for spikes.} Architectures for spike-time
prediction and for neural simulation have been explored with convolutional
networks, recurrent networks \cite{hochreiter1997long} and transformers
\cite{vaswani2017attention}. Deep networks have also been criticised for
requiring more data than electrophysiological datasets typically provide, a
critique that is partly addressed by structured or attention-based models
\cite{brunton2016why}. Our setup makes the data/parameter trade-off explicit by
reporting parameter counts alongside performance.

\textbf{Uncertainty.} Deep ensembles
\citep{lakshminarayanan2017ensembles} quantify epistemic uncertainty without
retraining, and split-conformal regression \cite{lei2018distribution} provides
finite-sample coverage guarantees under exchangeability. Exchangeability fails
exactly where the interesting science is: out of distribution. We therefore
report both, and we treat coverage \emph{degradation} as a diagnostic of
distribution shift rather than as a failure of the method.

\textbf{Closed-loop design.} Optimising behaviour by model-in-the-loop
optimisation is well established in reinforcement learning and in control. In
neuroscience, the closest practical analogue is closed-loop identification and
adaptation of neural dynamics. Our contribution is not to invent that loop but
to show that a model trained purely for prediction is a usable design surrogate,
and to quantify the design gain.

\section{The benchmark}
\label{sec:bench}

\subsection{Forward model}
\label{sec:sim}

Each virtual neuron is an adaptive exponential integrate-and-fire (AdEx) neuron
\cite{brette2005adex} with membrane time constant $\tau_m = C/g_L$ and adaptation
current $w$ obeying
\begin{align}
C\dot V &= -g_L(V-E_L) + g_L\Delta_T\exp\!\big(\tfrac{V-V_T}{\Delta_T}\big)
          - w + I_{\mathrm{bg}}(t) + I_{\mathrm{op}}(t), \\
\tau_w\dot w &= a(V - E_L) - w, \qquad V\!\ge\!V_{\mathrm{spk}}:\;
          V\!\leftarrow\!V_r,\; w\!\leftarrow\! w + b,
\end{align}
integrated with forward Euler at $\Delta t = 0.1$~ms.

The opsin is a two-state channelrhodopsin model. Opening is driven by a
Hill-shaped function of light intensity and deactivation is voltage-dependent,
so that the photocurrent is terminated during spikes:
\begin{equation}
\dot m = \frac{m_\infty(L) - m}{\tau_m(V)}, \qquad
\dot h = \frac{h_\infty(L) - h}{\tau_h(L)}, \qquad
I_{\mathrm{op}} = g_{\mathrm{opsin}}\, m\, h\, \sigma(V)\, (E_{\mathrm{ops}} - V),
\end{equation}
with $m_\infty(L) = L^{n}/(L^{n}+K^{n})$, $n=2$, $\sigma(V)$ a decreasing
logistic in voltage, and a slow desensitisation variable $h$ with
$\tau_h \approx 0.7$~s under light and $2.5$~s in darkness. The background
current is Ornstein--Uhlenbeck noise,
$\dot I = -I/\tau_{\mathrm{ou}} + \sigma\sqrt{2/\tau_{\mathrm{ou}}}\,dW$,
which supplies both a spontaneous firing rate and realistic trial-to-trial
variability.

Three cell types are simulated, differing in membrane capacitance, leak,
threshold, adaptation strength and opsin expression: a regular-spiking
pyramidal-like cell (\textsc{pc}), a fast-spiking interneuron-like cell
(\textsc{pv}), and an adapting interneuron-like cell (\textsc{som}). Baseline
drive for each type is calibrated by bisection so that spontaneous firing
rates are in a physiological range. Crucially, \emph{every} cell is a random
draw from a log-normal population in $C$, $g_L$, $E_L$, $V_T$, $\Delta_T$,
$a$, $b$, $\tau_w$, $g_{\mathrm{opsin}}$, $I_0$ and noise amplitude, with
coefficients of variation of $10$--$40\%$. A model that ignores cell identity
therefore cannot fit the data.

Figure~\ref{fig:sim} shows the three sanity checks we consider mandatory: a
threshold-like intensity--response relation with a graded supra-threshold
range; a following-failure relation in which spike fidelity per pulse falls
with frequency, more severely for the adapting cell type; and a realistic
response raster.

\paperfig{../results/figures/fig2_simulator_validation.png}{Forward-model validation. \textbf{Left:} mean evoked spike count as a
function of light intensity; the response is sub-threshold below
$\approx\!0.3$ model units and graded above it, as expected for a
conductance-based photocurrent with an activation threshold.
\textbf{Centre:} spikes per delivered pulse (following fidelity) against
pulse frequency at fixed intensity; fidelity falls with frequency, and the
strongly adapting pyramidal-like cell follows high-frequency trains worse than
the non-adapting fast-spiking cell. \textbf{Right:} example spike raster of
a \textsc{pv} cell to a 40~Hz, 5~ms train.}{fig:sim}

\subsection{Stimulation protocols and splits}

Each trial is a 1500~ms recording with light starting at 200~ms. Nine
protocol families are generated (Figure~\ref{fig:protos}): \emph{regular}
(log-uniform frequency 1--40~Hz trains), \emph{step} (single 10--800~ms pulse),
\emph{rand} (gamma-distributed inter-pulse intervals, CV 0.1--1.6),
\emph{doublet} (5--400~ms inter-pulse interval), \emph{burst} (burst trains with
2--8 pulses per burst), and four families reserved for OOD testing:
\emph{chirp} (exponential frequency sweep 2--60~Hz), \emph{sweep} (triangle
sweep), \emph{sparse} (very low rate, CV$>$1.8) and \emph{chirp\_hi}
(sweep to 70--120~Hz).

Splits are defined over both cells and protocols:
\begin{itemize}[leftmargin=*,itemsep=1pt]
\item \textbf{train}: 36,000 trials from the \nTrainFamilies{} training protocol
  families, drawn from 600 training cells;
\item \textbf{val}: 7,000 trials from 120 validation cells, used for early
  stopping and for conformal calibration;
\item \textbf{ID test}: 7,000 trials from the \emph{training} protocol
  distribution but \emph{unseen} cells (180 held-out cells);
\item \textbf{OOD protocol}: 5,000 trials from the four held-out families;
\item \textbf{OOD frequency}: 3,000 regular trains at 50--120~Hz;
\item \textbf{OOD intensity}: 3,000 protocols at intensities above the
  training range.
\end{itemize}
Holding out cells as well as protocols separates memorisation of a stimulus
from generalisation to a new cell, which is the regime that matters
experimentally.

\paperfig{../results/figures/fig1_protocol_examples.png}{One example stimulus from each of the nine protocol families. The
first five appear in training; chirp, sweep, sparse and chirp\_hi are held out
and constitute the OOD protocol split.}{fig:protos}

\paperfig{../results/figures/fig3_dataset_overview.png}{Benchmark composition.
\textbf{Left:} protocol family by split, showing that the OOD splits use
families never seen in training. \textbf{Centre:} distribution of evoked spike
counts per split (note the shift in the response distribution in the
OOD-protocol split). \textbf{Right:} training stimulus intensities per cell
type.}{fig:dataset}

\subsection{Features and targets}

Two prediction targets are defined per trial: the \emph{evoked spike count}
(spikes after stimulus onset) and the full 1~kHz \emph{spike raster}. The
latter is the primary task; the former is reported because it is interpretable
and is the quantity an experimenter designs for. Performance is measured with
(1)~count $\Rtwo$ and MAE, (2)~bin-level AUROC and AUPRC at 1~kHz, and (3)~PSTH
correlation, the per-trial Pearson correlation between the smoothed observed
raster and the predicted rate.

\section{Models}

All models share the same interface: input = 1~kHz light trace in
model-intensity units plus a cell descriptor (three cell-type indicators and
eleven normalised electrophysiological and opsin parameters); output = a
per-millisecond spike probability.

\paragraph{\textsc{glm} (classical encoding model).}
A cell-type-specific causal filter of $512$~ms followed by a pointwise
non-linearity, plus a linear contribution of the cell parameters,
\begin{equation}
\text{logit}(t) = \sum_{\tau} K_{\text{type}(\tau)}(t)\,\theta_\tau
                  + \beta_{\text{type}} + \bm{w}^\top \bm{z}_{\text{cell}} .
\end{equation}
This is the standard spike encoding model and is the baseline a deep model must
beat \cite{pillow2008spatiotemporal}.

\paragraph{\textsc{cnn} (causal dilated CNN).}
A strided stem aggregates light into 8~ms frames and injects the cell
descriptor as constant channels; seven residual dilated blocks
(dilations $2^{0..6}$, kernel 5, position-wise normalisation) expand the
effective receptive field beyond the trial; logits are upsampled back to
1~kHz.

\paragraph{\textsc{lstm} (convolutional LSTM).}
The same strided stem, followed by a single-layer LSTM with 96 hidden units,
upsampled to 1~kHz \cite{hochreiter1997long}.

\paragraph{\former{} (causal transformer).}
Light is embedded in $32$~ms patches with a learned within-patch positional
encoding; a cell-dependent FiLM modulation \cite{perez2018film} injects cell
identity; a causal transformer encoder \cite{vaswani2017attention} propagates
context; a linear head decodes each patch into per-millisecond logits.

\paragraph{\textsc{hgb} (gradient-boosted trees).}
Histogram gradient boosting \cite{friedman2001greedy} on handcrafted
features: cell identity and parameters, protocol family indicators,
intensity, pulse number, duration, duty cycle, frequency, width, coefficient of
variation, inter-burst frequency, and light dose. This model represents the
practically dominant alternative in the field, where one has hundreds of trials
rather than tens of thousands.

\paragraph{Causality.}
Every model is causal up to its analysis frame: the GLM is strictly causal
within its $512$~ms filter; the CNN and LSTM use 8~ms frames; \former{} uses
32~ms patches. Because a stimulation protocol is specified in advance, no
information about the \emph{neuron} beyond time $t$ is ever used, and the
stimulus lookahead is bounded by one frame. We enforce this with a unit test
that perturbs the light trace after a given time step and asserts that all
earlier predictions are unchanged; the test also guards the normalisation
layers, which is why position-wise (not group) normalisation is used in the CNN.

\section{Results}

\subsection{In-distribution prediction}

Table~\ref{tab:main} summarises in-distribution performance on held-out cells.
The classical GLM already reaches $\RtwoGlmId$ on evoked spike counts and
AUPRC $=\auPrcGlmId$ at 1~kHz, confirming that much of the stimulus--response
mapping is a time-invariant linear filter plus a non-linearity. All three deep
models improve on this, with \former{} achieving $\RtwoTransformerId$ and
PSTH correlation $=\psthTransformerId$, i.e. millisecond-accurate temporal
predictions, not merely accurate totals. Gradient-boosted trees on engineered
features reach $\RtwoHgbId$, which is respectable but does not dominate: at
$\sim\!0.2$~M features worth of trial count, the trees can capture nonlinear
interactions of summary statistics but not the fine temporal structure that
distinguishes a 40~Hz train from a burst pattern with the same pulse count.
Bootstrap 95\% confidence intervals over trials and seed-to-seed standard
deviations are reported in Table~\ref{tab:main}; across seeds the models are
highly stable ($\sigma_{seed} < 0.02$ $\Rtwo$).

\paperfig{../results/figures/fig5_psth_examples.png}{Mean PSTH across
held-out trials per cell type (black) against the \former{} prediction
(orange). The prediction tracks the observed mean, including the onset
transient and the decay after stimulus offset, while smoothing
trial-to-trial variability, as a conditional mean should.}{fig:psth}

\begin{table}[t]
  \centering
  % BEGIN AUTO MAIN TABLE
  \begin{tabular}{llccccc}
  \toprule
  Model & Params & $R^2$ (95\% CI) & $\sigma_{seed}$ & MAE & $\rho$ & Bin AUPRC \\
  \midrule
  GLM & $\nParamsGlm$ & $\RtwoGlmId$ ($\ciGlmLo$--$\ciGlmHi$) & $\seedSdGlm$ & $\maeGlmId$ & $\psthGlmId$ & $\auPrcGlmId$ \\
  CNN & $\nParamsCnn$ & $\RtwoCnnId$ ($\ciCnnLo$--$\ciCnnHi$) & $\seedSdCnn$ & $\maeCnnId$ & $\psthCnnId$ & $\auPrcCnnId$ \\
  LSTM & $\nParamsLstm$ & $\RtwoLstmId$ ($\ciLstmLo$--$\ciLstmHi$) & $\seedSdLstm$ & $\maeLstmId$ & $\psthLstmId$ & $\auPrcLstmId$ \\
  OptoFormer & $\nParamsTransformer$ & $\RtwoTransformerId$ ($\ciTransformerLo$--$\ciTransformerHi$) & $\seedSdTransformer$ & $\maeTransformerId$ & $\psthTransformerId$ & $\auPrcTransformerId$ \\
  HGB (trees) & --- & $\RtwoHgbId$ & --- & $\maeHgbId$ & --- & --- \\
  \bottomrule
  \end{tabular}
  % END AUTO MAIN TABLE
  \caption{In-distribution performance on held-out cells (7,000 trials).
  $\Rtwo$ and MAE refer to the evoked spike count; AUPRC is computed over all
  1~kHz bins (base rate $0.008$). Parenthesised values are 95\% bootstrap
  confidence intervals over trials. $\sigma_{seed}$ is the standard deviation of
  $\Rtwo$ across \nSeeds{} training seeds.}
  \label{tab:main}
\end{table}

\subsection{Out-of-distribution: the split that matters}

Table~\ref{tab:ood} is the central result. Gradient-boosted trees extrapolate
to unseen frequencies ($\RtwoHgbFreq$) and unseen intensities ($\RtwoHgbInt$)
about as well as they do in distribution, which reflects the fact that
frequency and intensity appear as smooth, monotone features. On entirely
unseen \emph{protocol families} the same model collapses to
$\RtwoHgbProto$: with no family indicator to key on and a train/regular
boundary it has never seen, its predictions are worse than the mean predictor.
This is the failure mode that a random-split benchmark cannot detect, and it
is exactly the failure mode that an experimenter designing a new stimulation
paradigm would encounter.

The sequence models, by contrast, degrade gracefully on the same split
($\RtwoTransformerProto$ for \former{}), because they read the stimulus
directly rather than through summary statistics. The ordering is stable across
all three OOD regimes and the paired per-trial comparisons are significant
(Section~\ref{sec:stats}).

\paperfig{../results/figures/fig4_main_results.png}{Main comparison.
\textbf{Left:} $\Rtwo$ of the evoked spike count per model and test regime.
\textbf{Centre:} bin-level AUPRC at 1~kHz, with the feature-based tree
baseline overlaid as a dashed line. \textbf{Right:} per-trial PSTH
correlation. Error bars are standard deviations across training seeds.}{fig:main}

\begin{table}[t]
  \centering
  % BEGIN AUTO OOD TABLE
  \begin{tabular}{lcccc}
  \toprule
  Model & ID & OOD protocol & OOD freq. & OOD intensity \\
  \midrule
  GLM & $\RtwoGlmId$ & $\RtwoGlmProto$ & $\RtwoGlmFreq$ & $\RtwoGlmInt$ \\
  CNN & $\RtwoCnnId$ & $\RtwoCnnProto$ & $\RtwoCnnFreq$ & $\RtwoCnnInt$ \\
  LSTM & $\RtwoLstmId$ & $\RtwoLstmProto$ & $\RtwoLstmFreq$ & $\RtwoLstmInt$ \\
  OptoFormer & $\RtwoTransformerId$ & $\RtwoTransformerProto$ & $\RtwoTransformerFreq$ & $\RtwoTransformerInt$ \\
  HGB (trees) & $\RtwoHgbId$ & $\RtwoHgbProto$ & $\RtwoHgbFreq$ & $\RtwoHgbInt$ \\
  \bottomrule
  \end{tabular}
  % END AUTO OOD TABLE
  \caption{$\Rtwo$ of evoked spike count for each model across the four test
  regimes. The tree baseline extrapolates in frequency and intensity but fails
  on unseen protocol families; sequence models degrade gracefully.}
  \label{tab:ood}
\end{table}

\noindent
One caveat deserves emphasis, because it affects how the negative values should
be read: the OOD-protocol split also shifts the \emph{response} distribution,
since chirp, sweep and sparse stimuli evoke fewer spikes on average than the
training families (Figure~\ref{fig:dataset}, centre). A model that predicts
near the training mean therefore incurs a large bias penalty, which is part of
what $\Rtwo < 0$ measures. The informative comparison is the \emph{relative}
ordering of models on the same split: models that read the stimulus directly
retain most of their accuracy, while the tree baseline does not.

\subsection{Predictive uncertainty detects the shift}

Deep ensembles over \nSeeds{} seeds plus split-conformal prediction give a
distribution-free interval at nominal $90\%$ coverage, calibrated on the
validation split. Coverage on the ID test is close to nominal
($\covTransformerId$ for \former{}), while coverage degrades on the OOD
splits, and the mean ensemble standard deviation increases
($\sdTransformerId \rightarrow \sdTransformerProto$ spikes per trial). The
conformal interval is therefore not only honest in-distribution, it also
\emph{warns}: a practitioner running protocol screening would see wide
intervals and high ensemble disagreement exactly when the model is being
extrapolated, which is the signal to fall back on a simpler protocol.

\subsection{Few-shot transfer to a held-out cell type}

We pretrained a model on two cell types and asked how much labelled data a
third type needs. Zero-shot performance after removing that type from training
is $\trZeroShotRtwo$; fine-tuning on 50 labelled cells of the held-out type
gives $\trFewShotNRtwo$, versus $\trScratchNRtwo$ when the same 50 cells are
used for training from scratch and $\trOracleRtwo$ for the oracle that saw the
type during pretraining. Fine-tuning the pretrained model on the same 50 cells
as the oracle ($\trFinetuneOracleNRtwo$) isolates the value of pretraining on
other types from the value of the extra data. The pattern---pretraining helps
most exactly when data are scarce---is the MAML-style intuition
\cite{finn2017maml}, here obtained without any meta-learning machinery.

\subsection{The model designs stimulation}

We used the surrogate (ensemble mean of \former{}) as an objective inside a
differential-evolution search over a ten-segment frequency schedule plus a
global intensity, targeting a prescribed evoked spike count
$\optTarget$ spikes per trial. Three search strategies were compared (learned
schedule, surrogate-optimised constant frequency, random search) plus an oracle
constant-frequency protocol scanned directly in the simulator. Every candidate
was finally scored in the ground-truth simulator on held-out cells of the target
type, so the comparison measures design quality, not surrogate agreement.

Mean absolute error of the realised population spike count with respect to the
target is $\optErrOurs$ for schedule-based design, $\optErrRand$ for random
search over the same schedule space, $\optErrConst$ for the surrogate-optimised
constant frequency, and only $\optErrOracle$ for a constant frequency scanned
directly in the simulator. Three observations follow.

First, within the surrogate-guided searches, protocol expressiveness matters
more than the search algorithm: the flexible schedule improves on the
constant-frequency restriction ($\optErrOurs$ versus $\optErrConst$) while
delivering less light ($\optDoseOurs$ versus $\optDoseConst$ dose units), and
random search, which uses no gradient information, matches differential
evolution ($\optErrRand$).

Second, the gap between every surrogate-guided design and the
simulator-scanned oracle quantifies the \emph{optimizer's curse}. The search
concentrates on designs where the surrogate is most confident, so any residual
bias the surrogate carries is amplified exactly there: a predictor that scores
$\RtwoTransformerId$ on held-out trials still designs protocols that miss the
target by more than an order of magnitude more than a direct scan. Closing
this gap requires ground-truth evaluations inside the loop, i.e.\ active
learning rather than pure model-based optimisation.

Third, no method controls single-cell outcomes: only $\optWithinOurs$ of
individual cells fall within $\pm3$ spikes of the target for the best design,
and the simulator-scanned oracle is no better ($\optWithinOracle$). Cell-to-cell
variability, not design quality, dominates single-cell agreement, which is why
the target is phrased as a population mean---the quantity these methods do
control ($\optErrOracle$ for the oracle).

\subsection{Interpretability recovers biophysics}

Gradient saliency shows where each model reads the stimulus. For the CNN,
mean $|\partial\,\text{logit}/\partial L|$ inside pulse windows exceeds
off-pulse saliency by $\salEnrichCnn$-fold, and for the GLM by
$\salEnrichGlm$; the recurrent models weight in- and off-pulse bins about
equally ($\salEnrichLstm$, $\salEnrichTransformer$), consistent with
sensitivity to extended light history rather than to instantaneous drive
alone. All models weight the beginning of the trial heavily: for \former{},
mean saliency in the first 500~ms after onset is $\salEarlyTransformer$
times that of the second half (Figure~\ref{fig:interp}, left). Crucially,
the negative control works: shuffling the light trace across trials destroys
prediction (AUPRC falls to $\shuffleAuPrcTransformer$, at the base rate of
$0.008$), demonstrating that the models use the stimulus rather than trial
identity or cell identity alone.

The learned GLM filters are dominated by a single short-latency component:
the peak of $|K|$ lies $\glmPeakPc$~ms after the most recent light sample for
pyramidal-like cells ($\glmPeakPv$~ms PV-like, $\glmPeakSom$~ms SOM-like),
with near-zero weight at longer lags (Figure~\ref{fig:interp}, centre) ---
the fast onset of the two-state ChR2 photocurrent used by the forward model.
Because the GLM kernel is a literal causal filter, its shape can be compared
directly against the photocurrent time course used by the forward model,
providing an independent check that the model has identified the right
mechanism rather than a correlated shortcut.

\former{}'s mean attention matrix (Figure~\ref{fig:interp}, right) shows no
sharply local pattern: only $\attnLocal$ of the attention mass falls within
three patches of the diagonal, against $\attnLocalUniform$ under uniform
causal attention. Late queries spread attention across the whole causal
window, with a mild preference for mid-stimulus patches --- the attention
picture thus agrees with the saliency picture: the deep models integrate
over extended stimulus history instead of reacting only to the
instantaneous light level.

\paperfig{../results/figures/fig6_uncertainty.png}{Uncertainty under distribution shift.
\textbf{Left:} empirical coverage of 90\% split-conformal intervals, calibrated
on the validation split; coverage degrades as the test regime moves out of
distribution. \textbf{Right:} mean ensemble standard deviation of the
predicted spike count, which grows in the same regimes --- the model signals
when it is extrapolating.}{fig:uq}

\paperfig{../results/figures/fig8_transfer.png}{Cross-cell-type transfer.
Prediction accuracy on the held-out cell type under five conditions: an oracle
that saw the type during pretraining, zero-shot transfer, few-shot
fine-tuning of the pretrained model on 50 cells of the held-out type,
training from scratch on those same 50 cells, and fine-tuning of the oracle on
them.}{fig:transfer}

\paperfig{../results/figures/fig9_protocol_optimization.png}{Model-in-the-loop stimulation design, per cell type.
\textbf{Left:} frequency schedules produced by surrogate-driven differential
evolution, by surrogate-optimised constant frequency, and by random search.
\textbf{Centre:} realised spike counts of the designed protocols when run in
the ground-truth simulator on held-out cells. \textbf{Right:} the
corresponding light protocols.}{fig:opt}

\paperfig{../results/figures/fig7_interpretability.png}{Interpretability.
\textbf{Left:} mean gradient saliency as a function of time after stimulus
onset: \former{} peaks at onset and decays, the CNN and LSTM oscillate with
the stimulation rhythm, and the GLM is flat. \textbf{Centre:} learned GLM
impulse responses by cell type: a single short-latency peak and near-zero
weight at longer lags. \textbf{Right:} mean causal attention between
patches; mass is spread across the causal window rather than concentrated
near the diagonal.}{fig:interp}

\section{Statistical analysis}
\label{sec:stats}

All comparisons use the same held-out cells. Confidence intervals for $\Rtwo$,
MAE and Spearman correlation are bootstrap percentile intervals over trials
($1{,}000$ resamples). Model comparisons use paired Wilcoxon signed-rank tests
on per-trial absolute error with Holm correction across comparisons. Seed
variability is reported as the standard deviation across \nSeeds{} training
seeds. Absolute effect sizes are reported alongside $p$-values, in keeping
with the observation that with thousands of trials, statistical significance
can accompany a practically trivial difference---so we emphasise the $\Rtwo$
and MAE differences themselves.

\section{Limitations and threats to validity}
\label{sec:limits}

\begin{enumerate}[leftmargin=*,itemsep=2pt]
\item \textbf{Simulation.} The forward model, however carefully built, is not
      cortex. A model trained here has learned the simulator's stimulus--response
      map. The transferable claim we make is about the \emph{methodology} and
      about the failure modes, not about a specific prediction for a real
      neuron. Real-data validation should re-fit the light-to-spike map on
      patch-clamp or two-photon recordings and then apply the same OOD
      protocol to test whether the degradation pattern reported here
      reproduces.
\item \textbf{Single-neuron, single-photon class.} We model direct
      opsin-driven excitation of three cell types. Recurrent network
      interactions, opsin inhibition, and multi-wavelength designs are out of
      scope, though the simulator can be extended to them and the code makes
      this straightforward.
\item \textbf{Model scale.} The models here are small ($\sim\!10^{4}$--$10^{5}$
      parameters) because the benchmark is generated, not because small is
      optimal. We expect performance to improve with scale and with more
      protocol diversity; we report the data/parameter trade-off explicitly
      rather than tuning per model to a fixed compute budget that would bias
      the comparison.
\item \textbf{Frame-level stimulus lookahead.} CNN and LSTM operate on 8~ms
      frames and \former{} on 32~ms patches, so predictions may use stimulus
      information from within the current frame. Since protocols are specified
      in advance, this does not leak neuronal information; it does mean the
      models are not usable for reacting to a stimulus that is not known in
      advance, which matters only for truly online closed-loop use.
\end{enumerate}

\section{Conclusion}

Optogenetics is limited less by our ability to deliver light than by our
ability to predict what the light will do. We built a benchmark on which that
question has a checkable answer, and showed that (i)~deep sequence models
outperform both a classical encoding model and gradient-boosted trees on
held-out cells, (ii) the classical tree approach fails precisely where novel
stimulation paradigms live while sequence models do not, (iii)~ensembles and
conformal intervals expose that failure rather than hiding it, (iv)~a few dozen
cells adapt a model to a new cell type, and (v)~a model trained only for
prediction is a usable, if biased, design surrogate whose bias is revealed by
scoring its designs in the simulator. The practical
consequence is a shorter experimental loop: screen, design, validate, repeat.

\section*{Data and code availability}

All code, the data generator, configuration files, tests and result tables are
released with this manuscript in the accompanying repository at
\url{https://github.com/basitali08/optonet}.
The dataset is regenerated deterministically from a single seed with
\texttt{python scripts/01\_generate\_data.py}.

\section*{Acknowledgements}

We thank the open-source communities behind PyTorch, NumPy, SciPy,
scikit-learn and Matplotlib.

\bibliographystyle{plainnat}
\bibliography{refs}

\end{document}
